\documentclass[11pt,a4paper]{article}
\usepackage[T1]{fontenc}
\usepackage[utf8]{inputenc}
\usepackage{amsmath,amssymb,amsthm}
\usepackage[margin=1in]{geometry}
\usepackage[colorlinks=true,linkcolor=blue,urlcolor=blue]{hyperref}
\theoremstyle{plain}
\newtheorem{theorem}{Theorem}
\newtheorem{lemma}[theorem]{Lemma}
\newtheorem{proposition}[theorem]{Proposition}
\newtheorem{corollary}[theorem]{Corollary}
\theoremstyle{definition}
\newtheorem{remark}[theorem]{Remark}
% A quoted conjecture can be a single hundred-term formula whose terms are thirty-digit
% numbers. TeX cannot hyphenate those, so without a little stretch every such quote runs off
% the page; the linked-a-file papers are all of that shape.
\setlength{\emergencystretch}{3em}

\title{The empirical recurrence for OEIS A233846, proved by transfer matrix}
\author{Adrian Perez Fontelles\\ \small Independent researcher}
\date{14 September 2026}
\begin{document}
\maketitle

\begin{abstract}
OEIS A233846 counts arrays of width $3$ over $\{0,1,2,3\}$ containing no run of $3$ consecutive
cells, taken along a named direction of the grid, whose entries are strictly increasing. The entry carries
an empirical recurrence of order $56$ contributed by R.~H.~Hardin. It is true, and it is
decidable rather than empirical: every such run lies inside $3$ consecutive array rows, so
the count is a walk count on $S=4161$ states.
 Throughout, \emph{row} means a line of the array in the direction the walk grows: this entry fixes the first dimension and grows the second, so the array is read transposed, and a step appends what the entry's own name writes as a column.\end{abstract}

\noindent\small 2020 Mathematics Subject Classification. 05A15, 68R15, 15A18.\normalsize

\section{The conjecture}

OEIS A233846 is ``Number of (1+2)X(n+2) 0..3 arrays with no increasing sequence of length 3 horizontally or antidiagonally downwards.''. It has offset $1$ and begins
\[
201664, 9828513, 480071296, 23453646656, 1144779904896, 55884483506105, 2727732140106016, 133143607148769360,\ \dots
\]
The entry states, as a conjecture contributed by its author and never marked settled:
\begin{quote}
Empirical recurrence of order 56 (see link above)

Empirical: a(n) = 60*a(n-1) - a(n-2) - 33104*a(n-3) + 212786*a(n-4) + 6488828*a(n-5) - 60862818*a(n-6) - 510400764*a(n-7) + 6669104935*a(n-8) + 10453760096*a(n-9) - 320921068017*a(n-10) + 317555763536*a(n-11) + 8196969426516*a(n-12) - 20162187878008*a(n-13) - 119384123556824*a(n-14) + 451862708541684*a(n-15) + 956745558062738*a(n-16) - 5629816801581596*a(n-17) - 2906458723857302*a(n-18) + 42652858845354660*a(n-19) - 15645613778200832*a(n-20) - 200231786265874500*a(n-21) + 189444976544641624*a(n-22) + 580878964472074828*a(n-23) - 824824276609672294*a(n-24) - 1040105754911889916*a(n-25) + 2062576702609464426*a(n-26) + 1083579763768407652*a(n-27) - 3323309570266011964*a(n-28) - 433815807412070012*a(n-29) + 3613479839578962524*a(n-30) - 430180494508749380*a(n-31) - 2691493521755884737*a(n-32) + 806975581039469336*a(n-33) + 1361866624557434527*a(n-34) - 617560391219808948*a(n-35) - 451367287392865314*a(n-36) + 279903967985979320*a(n-37) + 89240983439929962*a(n-38) - 79588772821592320*a(n-39) - 7468431410353569*a(n-40) + 14107899140663340*a(n-41) - 630894458822057*a(n-42) - 1507628117908676*a(n-43) + 211646084396264*a(n-44) + 92130011750340*a(n-45) - 20967063414020*a(n-46) - 2752736919552*a(n-47) + 1029479499808*a(n-48) + 10153795896*a(n-49) - 25179423448*a(n-50) + 1377418096*a(n-51) + 253463280*a(n-52) - 24890944*a(n-53) - 602592*a(n-54) + 130048*a(n-55) - 3600*a(n-56) for n>58
\end{quote}
As of the ``Last modified'' line on the live entry (Jul 23 2025 08:12:58, revision 8) this is still
recorded as empirical, and nothing on the entry records it as proved.

\section{What is being counted}

A RUN is $3$ consecutive cells taken along one of the directions
\[
\begin{array}{l}
\text{down an antidiagonal, ne to sw, read the other way} \\
\text{down a column} \\
\end{array}
\]
all of them inside the array. The array is counted when no run has its entries strictly increasing. The entry writes the array with $n$ as the number of COLUMNS; the walk is run down the transpose, so the directions above are already the transposed ones.


One direction above is marked ``read the other way''. Transposing an array exchanges a row with
a column and carries each diagonal direction to itself only up to sense: an antidiagonal read
from the northeast becomes one read from the southwest. For a condition built from EQUALITY
that would not matter, but ``strictly increasing'' read backwards is ``strictly decreasing'',
so the sense is carried explicitly rather than normalised away. The
difference is not subtle: normalising it away makes the model disagree with the entry's
published terms from the very first one, which is how the point was found.

The entry carries no renaming clause, and it could not: ``strictly increasing'' compares
letters by size, so it is not preserved by a permutation of the alphabet. The count is taken
over the alphabet the entry names.

\section{Every run lies in $3$ consecutive rows}

\begin{lemma}\label{lem:win}
Each direction above changes the row index by $0$ or $1$ from one cell of a run to the next, so
a run of $3$ cells lies inside $3$ consecutive array rows, and a run whose last cell is in
row $i$ lies inside rows $i-2,\dots,i$.
\end{lemma}

\begin{proof}
Immediate from the row steps of the four directions, which are $0$ for a row and $1$ for the
other three.
\end{proof}

So the state is the window of the last $2$ array rows, a row not yet written carried as
$\bot$; a step appends a row and rejects it exactly when some run ends in it or lies inside it.
A run needing a row that is $\bot$ does not exist, which is the boundary rule and is read off
the window rather than assumed. The walk starts at the all-$\bot$ window, so a walk of $R$ steps
is an array of $R$ rows. With $M$ the adjacency matrix of that digraph,
\[
a(n)\;=\;\iota^{\!\top}M^{\,j}\tau ,\qquad S=4161,\quad\tau=\mathbf{1} .
\]
The walk index $j$ and the entry's own $n$ differ by a constant, read off by matching the
published terms rather than assumed. In particular $a$ is $C$-finite.

\section{The proof}

\begin{lemma}\label{lem:crit}
Let $q(t)=t^{56}-\sum_i c_it^{56-i}$ be the characteristic polynomial of the
conjectured recurrence. The residual $a(n)-\sum_ic_ia(n-i)$ equals
$\iota^{\!\top}M^{\,j}q(M)\tau$ for the corresponding walk index $j$, so the recurrence
holds from some point on if and only if $\iota^{\!\top}M^{j}q(M)\tau=0$ for all large $j$,
and it is enough to examine $j<S$.
\end{lemma}

\begin{proof}
Factor $M^{j}$ out of $M^{j+56}-\sum_ic_iM^{j+56-i}$. For the bound, $M^{S}$
is an integer combination of $I,M,\dots,M^{S-1}$ by Cayley--Hamilton, so the numbers
$u_j=\iota^{\!\top}M^{j}q(M)\tau$ obey the monic recurrence given by the characteristic
polynomial of $M$; $S$ consecutive zeros force every later one, and the last nonzero $u_j$
pins the threshold exactly.
\end{proof}

\begin{theorem}
\[
a(n)\;=\;60\,a(n-1) - a(n-2) - 33104\,a(n-3) + 212786\,a(n-4) + 6488828\,a(n-5) - 60862818\,a(n-6) - 510400764\,a(n-7) + 6669104935\,a(n-8) + 10453760096\,a(n-9) - 320921068017\,a(n-10) + 317555763536\,a(n-11) + 8196969426516\,a(n-12) - 20162187878008\,a(n-13) - 119384123556824\,a(n-14) + 451862708541684\,a(n-15) + 956745558062738\,a(n-16) - 5629816801581596\,a(n-17) - 2906458723857302\,a(n-18) + 42652858845354660\,a(n-19) - 15645613778200832\,a(n-20) - 200231786265874500\,a(n-21) + 189444976544641624\,a(n-22) + 580878964472074828\,a(n-23) - 824824276609672294\,a(n-24) - 1040105754911889916\,a(n-25) + 2062576702609464426\,a(n-26) + 1083579763768407652\,a(n-27) - 3323309570266011964\,a(n-28) - 433815807412070012\,a(n-29) + 3613479839578962524\,a(n-30) - 430180494508749380\,a(n-31) - 2691493521755884737\,a(n-32) + 806975581039469336\,a(n-33) + 1361866624557434527\,a(n-34) - 617560391219808948\,a(n-35) - 451367287392865314\,a(n-36) + 279903967985979320\,a(n-37) + 89240983439929962\,a(n-38) - 79588772821592320\,a(n-39) - 7468431410353569\,a(n-40) + 14107899140663340\,a(n-41) - 630894458822057\,a(n-42) - 1507628117908676\,a(n-43) + 211646084396264\,a(n-44) + 92130011750340\,a(n-45) - 20967063414020\,a(n-46) - 2752736919552\,a(n-47) + 1029479499808\,a(n-48) + 10153795896\,a(n-49) - 25179423448\,a(n-50) + 1377418096\,a(n-51) + 253463280\,a(n-52) - 24890944\,a(n-53) - 602592\,a(n-54) + 130048\,a(n-55) - 3600\,a(n-56)
\]
for every $n>58$.
\end{theorem}

\begin{proof}
The vector $w=q(M)\tau$ was formed by $56$ matrix--vector products in exact integer
arithmetic and $\iota^{\!\top}M^{j}w$ evaluated for $j=0,1,\dots$, again exactly; those
integers vanish from the index corresponding to $n=58+1$ onwards and the run continued
until the working vector was identically zero.
\end{proof}

\section{Verification}

Four checks, each independent of the algebra above.

First, the model reproduces all $12$ terms the entry publishes, exactly, in integer
arithmetic.

Second, the count was recomputed for the smallest arrays straight from the definition: fill the
cells in row major order, in the orientation the entry's own name writes, and reject as soon as
a completed run is strictly increasing. No window, no transposition and no transfer matrix are involved, and
the published terms came back.

Third, the conjectured recurrence was evaluated directly on the published terms, with no
matrices involved, and holds wherever the proved range and the data overlap.

Fourth, the annihilation test of Lemma~\ref{lem:crit} was carried out in exact integer
arithmetic on the whole vector rather than on a sampled prefix.

\begin{thebibliography}{9}
\bibitem{oeis} The OEIS Foundation, \emph{The On-Line Encyclopedia of Integer
Sequences}, \url{https://oeis.org}, 2026. Sequence A233846.
\bibitem{stanley} R.~P.~Stanley, \emph{Enumerative Combinatorics}, Volume 1, 2nd ed.,
Cambridge University Press, 2011. (Transfer-matrix method, Section 4.7.)
\bibitem{horn} R.~A.~Horn and C.~R.~Johnson, \emph{Matrix Analysis}, 2nd ed., Cambridge
University Press, 2013.
\end{thebibliography}

\end{document}
